% ==============================================================================
% Chapter 3 — Data and Methods
% ==============================================================================

\chapter{Data and Methods}
\label{ch:methods}

\todo[inline]{Write data and methods chapter.}

\section{Emission Input Data}
\label{sec:data_emissions}

\subsection{CAMS-REG RI-URBANS Inventory}
\label{sec:data_cams}

The \cams{} emission inventory used in this study is the RI-URBANS version of
CAMS-REG-v6.1 \citep{Kuenen_ST15_2025}, with a base year of 2019 and a
horizontal resolution of approximately $\sim$\qtyproduct{6 x 6}{\km}. The
dataset covers 13 GNFR emission sectors and the following pollutants: NO$_x$,
PM\textsubscript{2.5}, PM\textsubscript{10}, NMVOC, SO\textsubscript{2},
NH\textsubscript{3}, CO, and total particle number (TPN).

\todo[inline]{Describe file format (NetCDF), units (Tg/year), grid definition,
and how data were obtained from the TNO FTP repository.}

\subsection{UFP Number Inventory}
\label{sec:data_ufp}

\todo[inline]{Describe TPN emissions data: size-speciated UFP, sectors covered,
data availability.}

\section{Spatial Proxy Datasets}
\label{sec:data_proxies}

\begin{table}[h]
  \centering
  \caption{Spatial proxy datasets used for emission downscaling.}
  \label{tab:proxies}
  \begin{tabular}{llll}
    \toprule
    \textbf{Proxy} & \textbf{Source} & \textbf{Resolution} & \textbf{GNFR sectors} \\
    \midrule
    Population density    & GHS-POP 2015 (JRC)   & 30 arc-sec & C, E \\
    Road network          & OpenStreetMap        & Vector     & F \\
    Traffic volumes       & OpenTransportMap     & Vector     & F (weighting) \\
    Industrial areas      & CORINE Land Cover 2018 & 100 m   & A, B, D \\
    Agricultural areas    & CORINE Land Cover 2018 & 100 m   & K, L \\
    Port areas            & CORINE Land Cover 2018 & 100 m   & G \\
    Airport areas         & CORINE Land Cover 2018 & 100 m   & H \\
    Railway network       & OpenStreetMap        & Vector     & I (railway) \\
    Point sources         & E-PRTR v17           & Point      & A, B, D, J \\
    Urban centre          & GHS-UCDB 2015        & Polygon    & F (weighting) \\
    \bottomrule
  \end{tabular}
\end{table}

\todo[inline]{Describe each proxy in detail: data source, preprocessing steps,
known limitations.}

\section{Study Domain Definition}
\label{sec:data_domain}

\todo[inline]{Define the model domains for Ioannina and Kozani: bounding box in
UTM coordinates, grid resolution (1x1 km), coordinate reference system.
Include a map figure showing both domains.}

\section{The UrbEm Downscaling Framework}
\label{sec:methods_urbem}

\todo[inline]{Give a concise description of the UrbEm methodology
\citep{Ramacher2021}: area/line/point source workflow, proxy normalisation,
mass conservation, output format for EPISODE-CityChem.}

\section{EPISODE-CityChem Simulation Setup}
\label{sec:methods_episode}

\todo[inline]{Describe model configuration: domain, meteorology input,
chemistry scheme, boundary conditions, simulation period.}

\section{Evaluation Methodology}
\label{sec:methods_eval}

\todo[inline]{Describe the observational datasets used for model evaluation
(monitoring stations, species, periods) and the statistical metrics
(RMSE, bias, correlation, FAC2).}
